\section{Three-dimensional applications}
\label{sec:three-dimensional-applications}

\subsection{Explicit closure for small critical data}

The construction below gives an explicit three-dimensional
\(\dot H^{1/2}\) realization of the critical-data mechanism introduced by
Fujita and Kato \cite{FujitaKato1964}; every constant used here is derived
in the displayed estimates.

Use the explicit homogeneous Sobolev constants
\[
 \mathsf S_3:=\mathsf S(3),
 \qquad
 C_{\rm src}:=\mathsf S_3^3=\frac1{\sqrt2\,\pi}.
\]
The same constants hold for vector fields.  Indeed, apply the scalar
inequality to the Euclidean length \(|F|\): one has
\(|\nabla|F||\le|\nabla F|\), while the difference-quotient representation of
\(\dot H^{1/2}\) and
\(\bigl||F(x)|-|F(y)|\bigr|\le|F(x)-F(y)|\) give
\(\norm{|F|}_{\dot H^{1/2}}\le\norm F_{\dot H^{1/2}}\).

\begin{lemma}[Critical action identity]
\label{lem:app-critical-action-identity}
Let \((u,p)\) be a smooth pressure-normalized solution on a time interval
\(I\subset[0,T_*)\), and set
\[
 \mathcal B(u):=\Pp\nabla\!\cdot(u\otimes u).
\]
Then, pointwise on \(I\),
\begin{equation}
 \nu\frac{\dd}{\dd t}\norm{u}_{\dot H^{1/2}}^2
 +\nu^2\norm{u}_{\dot H^{3/2}}^2
 +\norm{u_t}_{\dot H^{-1/2}}^2
 =\norm{\mathcal B(u)}_{\dot H^{-1/2}}^2.
 \label{eq:app-critical-action-identity}
\end{equation}
\end{lemma}

\begin{proof}
The projected equation is
\[
 u_t=-\nu\Lambda^2u-\mathcal B(u).
\]
Put
\[
 g:=\Lambda^{-1/2}\mathcal B(u),
 \qquad
 w:=\Lambda^{-1/2}u_t=-\nu\Lambda^{3/2}u-g.
\]
Pairing the projected equation with \(\Lambda u\) gives
\[
 \frac{\dd}{\dd t}\norm u_{\dot H^{1/2}}^2
 =-2\nu\norm u_{\dot H^{3/2}}^2
  -2\operatorname{Re}\ip g{\Lambda^{3/2}u}_2.
\]
Multiply this equality by \(\nu\) and expand
\(\norm w_2^2=\norm{\nu\Lambda^{3/2}u+g}_2^2\).
The cross terms cancel and give
\eqref{eq:app-critical-action-identity}.
\end{proof}

For \(m\in\N\), define the explicit homogeneous commutator constant
\begin{equation}
 \mathsf K_m^{\rm hom}:=
 3^m\sum_{\ell=1}^m\binom m\ell
 \mathsf S(\mathfrak p_{m,\ell})\mathsf S(\mathfrak q_{m,\ell}),
 \label{eq:app-explicit-homogeneous-commutator-constant}
\end{equation}
where \((\mathfrak p_{1,1},\mathfrak q_{1,1})=(3,6)\), and, for \(m\ge2\),
\[
 \vartheta_{m,\ell}:=\frac{\ell-1}{m-1},
 \qquad
 \frac1{\mathfrak p_{m,\ell}}:=\frac13-\frac{\vartheta_{m,\ell}}6,
 \qquad
 \frac1{\mathfrak q_{m,\ell}}:=\frac16+\frac{\vartheta_{m,\ell}}6.
\]
The corresponding estimate is
\begin{equation}
 \abs{\operatorname{Re}
 \ip{\Lambda^m\Pp\nabla\!\cdot(z\otimes z)}
     {\Lambda^mz}_{L^2}}
 \le
 \mathsf K_m^{\rm hom}
 \norm{\Lambda^{3/2}z}_2
 \norm{\Lambda^mz}_2
 \norm{\Lambda^{m+1}z}_2
 \label{eq:app-high-order-critical-commutator}
\end{equation}
for every \(z\in\SSigma(\R^3)\).

\begin{lemma}[Explicit homogeneous critical commutator]
\label{lem:app-critical-commutator}
For every \(m\in\N\), \(\mathsf K_m^{\rm hom}<\infty\), and
\eqref{eq:app-high-order-critical-commutator} holds with the displayed
constant.
Moreover, \eqref{eq:app-high-order-critical-commutator} extends to every
divergence-free
\(z\in H^{m+1}(\R^3)\cap\dot H^{3/2}(\R^3)\).
\end{lemma}

\begin{proof}
The multinomial identity gives the exact pairing
\[
 \ip{\Lambda^mf}{\Lambda^mg}_{L^2}
 =
 \sum_{|\alpha|=m}\frac{m!}{\alpha!}
 \ip{\partial^\alpha f}{\partial^\alpha g}_{L^2}.
\]
After differentiating \((z\cdot\nabla)z\), the undifferentiated transport
term has zero pairing by incompressibility.  A remaining term has the form
\[
 \ip{(\partial^\beta z\cdot\nabla)
          \partial^{\alpha-\beta}z}
      {\partial^\alpha z}_{L^2},
 \qquad
 |\alpha|=m,\quad 0<\beta\le\alpha.
\]
Put \(\ell=|\beta|\).  For \(m\ge2\), H\"older's inequality,
Lemma~\ref{lem:explicit-homogeneous-sobolev}, and Fourier log-convexity give
\begin{align*}
 &\norm{\partial^\beta z}_{\mathfrak p_{m,\ell}}
 \norm{\nabla\partial^{\alpha-\beta}z}_{\mathfrak q_{m,\ell}}\\
 &\quad\le
 \mathsf S(\mathfrak p_{m,\ell})
 \mathsf S(\mathfrak q_{m,\ell})
 \norm z_{\dot H^{3/2}}\norm z_{\dot H^{m+1}}.
\end{align*}
Indeed, the two differentiated Sobolev orders are the complementary
interpolants
\begin{align*}
 \ell+s(\mathfrak p_{m,\ell})
 &=(1-\vartheta_{m,\ell})\frac32
   +\vartheta_{m,\ell}(m+1),\\
 m-\ell+1+s(\mathfrak q_{m,\ell})
 &=\vartheta_{m,\ell}\frac32
   +(1-\vartheta_{m,\ell})(m+1).
\end{align*}
For \(m=1\), the same bound is the direct \(L^3\)-\(L^6\) estimate.
Finally,
\[
 \sum_{|\alpha|=m}\frac{m!}{\alpha!}=3^m,
 \qquad
 \sum_{\substack{\beta\le\alpha\\|\beta|=\ell}}
 \binom\alpha\beta=\binom m\ell.
\]
Summation gives exactly
\eqref{eq:app-explicit-homogeneous-commutator-constant} and proves
\eqref{eq:app-high-order-critical-commutator}.
Solenoidal Fourier truncation and convergence in
\(H^{m+1}\cap\dot H^{3/2}\) give the asserted extension.
\end{proof}

\begin{theorem}[Explicit small-critical-data global regularity]
\label{thm:app-small-critical-global}
Let \(\nu>0\) and \(a\in\SSigma(\R^3)\).
Set
\[
 a_{\rm c}:=\norm{\Lambda^{1/2}a}_2,
 \qquad
 \delta_0:=\nu^2-C_{\rm src}^2a_{\rm c}^2.
\]
If
\begin{equation}
 a_{\rm c}<\frac{\nu}{C_{\rm src}}=\sqrt2\,\pi\nu,
 \label{eq:app-critical-smallness}
\end{equation}
then \(a\) generates a unique global pressure-normalized smooth solution
\[
 (u,p)\in C^\infty([0,\infty)\times\R^3;\R^3\times\R).
\]
For every finite \(T>0\), the solution is the unique member of
\(\mathscr M_\nu^3(0,T;a)\).
For every \(t\ge0\),
\begin{equation}
 \nu\norm{\Lambda^{1/2}u(t)}_2^2
 {}+\delta_0\int_0^t\norm{\Lambda^{3/2}u(s)}_2^2\dd s
+\int_0^t\norm{u_t(s)}_{\dot H^{-1/2}}^2\dd s
 \le\nu a_{\rm c}^2,
 \qquad \delta_0>0.
 \label{eq:app-critical-stock-action}
\end{equation}
The critical nonlinear source and the first temporal response satisfy
\begin{align}
 \int_0^\infty
 \norm{\Pp\nabla\!\cdot(u\otimes u)(t)}_{\dot H^{-1/2}}^2\dd t
 &\le \frac{C_{\rm src}^2\nu a_{\rm c}^4}{\delta_0}
 =\frac{\nu a_{\rm c}^4}{2\pi^2\nu^2-a_{\rm c}^2},
 \label{eq:app-small-critical-nonlinear-bound}\\
 \nu^2\int_0^\infty\norm{\Lambda^{3/2}u(t)}_2^2\dd t
 +\int_0^\infty\norm{u_t(t)}_{\dot H^{-1/2}}^2\dd t
 &\le \nu a_{\rm c}^2+\frac{C_{\rm src}^2\nu a_{\rm c}^4}{\delta_0}.
 \label{eq:app-small-first-temporal-response}
\end{align}
For every \(m\in\N\), define
\[
 \Gamma_m^{\rm sm}
 :=\frac{(\mathsf K_m^{\rm hom})^2a_{\rm c}^2}{\delta_0}.
\]
Then
\begin{align}
 &\norm{\Lambda^mu(t)}_2^2
 +\nu\int_0^t\norm{\Lambda^{m+1}u(s)}_2^2\dd s
 \notag\\
 &\qquad\le
 \norm{\Lambda^ma}_2^2\exp(\Gamma_m^{\rm sm}).
 \label{eq:app-small-data-finite-order}
\end{align}
In particular,
\begin{equation}
 \sup_{t\ge0}\norm{u(t)}_{H^m}^2
 \le
 2^{m-1}\left[
 \norm a_2^2+
 \norm{\Lambda^ma}_2^2\exp(\Gamma_m^{\rm sm})
 \right].
 \label{eq:app-small-data-Hm}
\end{equation}
\end{theorem}

\begin{proof}
Iterate Theorem~\ref{thm:local-restart} and join overlapping intervals by
its at-most-one assertion for the shifted mild class.  The union is the
maximal shifted pressure-normalized \(H^3\) mild history on \([0,T_*)\), and
persistence retains \(H^\infty\) on
every compact subinterval.  Put
\[
 A(t):=\norm{\Lambda^{1/2}u(t)}_2,
 \qquad
 B(t):=\norm{\Lambda^{3/2}u(t)}_2.
\]
 Dual Sobolev, H\"older, and the unit Fourier-multiplier norm of \(\Pp\)
give the complete source-constant chain
\begin{align}
 \norm{\Pp\nabla\!\cdot(u\otimes u)}_{\dot H^{-1/2}}
 &\le \mathsf S_3\norm{(u\cdot\nabla)u}_{L^{3/2}}
 \notag\\
 &\le \mathsf S_3\norm u_{L^3}\norm{\nabla u}_{L^3}
 \notag\\
 &\le \mathsf S_3^3
      \norm{\Lambda^{1/2}u}_2
      \norm{\Lambda^{3/2}u}_2
 =C_{\rm src}A(t)B(t).
 \label{eq:app-critical-source-constant-chain}
\end{align}
Combining \eqref{eq:app-critical-source-constant-chain} with the exact
identity \eqref{eq:app-critical-action-identity} gives, on every classical
subinterval,
\[
 \nu\frac{\dd}{\dd t}A(t)^2
 +\bigl(\nu^2-C_{\rm src}^2A(t)^2\bigr)B(t)^2
 +\norm{u_t(t)}_{\dot H^{-1/2}}^2\le0.
\]
Continuity and a first-exit argument give \(A(t)\le a_{\rm c}\) throughout
\([0,T_*)\).  Integration and the definition of \(\delta_0\) prove
\eqref{eq:app-critical-stock-action}.  In particular,
\[
 \int_0^tB(s)^2\dd s\le\frac{\nu a_{\rm c}^2}{\delta_0}.
\]
For every \(t<T_*\), the same source estimate and the integrated identity
\eqref{eq:app-critical-action-identity} prove
\eqref{eq:app-small-critical-nonlinear-bound} and
\eqref{eq:app-small-first-temporal-response} with \(\infty\) replaced by
\(t\).  After the continuation argument below proves \(T_*=\infty\),
monotone convergence gives the stated forms.

At order \(m\), Lemma~\ref{lem:app-critical-commutator} and Young's
inequality yield
\[
 \frac{\dd}{\dd t}\norm{\Lambda^mu}_2^2
 +\nu\norm{\Lambda^{m+1}u}_2^2
 \le
 \frac{(\mathsf K_m^{\rm hom})^2}{\nu}
 B(t)^2\norm{\Lambda^mu}_2^2.
\]
Multiplication by the corresponding integrating factor and
\[
 \int_0^tB(s)^2\dd s\le\frac{\nu a_{\rm c}^2}{\delta_0}
\]
prove \eqref{eq:app-small-data-finite-order}.  The Fourier inequality
\[
 (1+\abs\xi^2)^m\le2^{m-1}(1+\abs\xi^{2m})
\]
and the kinetic-energy identity prove \eqref{eq:app-small-data-Hm}.

Let \(M_3\) be the square root of the right-hand side of
\eqref{eq:app-small-data-Hm} with \(m=3\), and set
\[
 \delta_{\rm crit}
 :=\min\{1,c_\nu(1+M_3)^{-2}\}>0.
\]
Every datum \(u(t_0)\), \(t_0<T_*\), has \(H^3\)-norm at most \(M_3\).
Theorem~\ref{thm:local-restart} restarts the solution for at least
\(\delta_{\rm crit}\) units of time.  If \(T_*<\infty\), choose
\(t_0<T_*\) with \(T_*-t_0<\delta_{\rm crit}/2\).  Both restrictions belong
to \(\mathscr M_\nu^3(t_0,b;u(t_0))\) for every \(b<T_*\), so uniqueness on
each such closed subinterval identifies them throughout \([t_0,T_*)\).
The restarted member then extends the original history beyond \(T_*\),
contradicting maximality.  Hence \(T_*=\infty\), and persistence
 gives the asserted smoothness.  The same local uniqueness on each finite
 interval proves the stated global uniqueness class.
\end{proof}

For \(m\ge1\), define
\[
 U_m^{\rm sm}
 :=
 2^{(m-1)/2}
 \left[
 \norm a_2^2+
 \norm{\Lambda^ma}_2^2\exp(\Gamma_m^{\rm sm})
 \right]^{1/2},
 \qquad
 U_0^{\rm sm}:=\norm a_2,
 \qquad
 U_{-1}^{\rm sm}:=U_0^{\rm sm}.
\]
For \(n\in\Nzero\) and \(m\ge-1\), define the finite-order recursion
\begin{align}
 \mathsf Y_{0,m}^{\rm sm}&:=U_m^{\rm sm},
 \label{eq:app-small-Y0}\\
 \mathsf Y_{n+1,m}^{\rm sm}
 &:=
 \nu\mathsf Y_{n,m+2}^{\rm sm}
 +A_{m+1}\sum_{a'=0}^{n}\binom n{a'}
 \mathsf Y_{a',m+2}^{\rm sm}
 \mathsf Y_{n-a',m+2}^{\rm sm}.
 \label{eq:app-small-Yrec}
\end{align}

\begin{proposition}[Explicit small-data pressure-complete rectangles]
\label{prop:app-small-data-rectangles}
Under the hypotheses of Theorem~\ref{thm:app-small-critical-global}, for
every \(n\in\Nzero\) and \(m\ge-1\),
\begin{equation}
 \sup_{t\ge0}\norm{V_n(t)}_{H^m}
 \le\mathsf Y_{n,m}^{\rm sm}.
 \label{eq:app-small-time-row}
\end{equation}
For every \(T<\infty\) and \(N,M\in\Nzero\),
\begin{align}
 \mathfrak R_{N,M}(u;T)
 &\le
 T\sum_{n=0}^{N}
 \left[
 (\mathsf Y_{n,M+1}^{\rm sm})^2
 +(\mathsf Y_{n+1,M-1}^{\rm sm})^2
 \right],
 \label{eq:app-small-rectangle}\\
 \mathfrak Q_{N,M}(p;T)
 &\le
 TA_M\sum_{n=0}^{N}\sum_{a'=0}^{n}\binom n{a'}
 \mathsf Y_{a',M+1}^{\rm sm}
 \mathsf Y_{n-a',M+1}^{\rm sm}.
 \label{eq:app-small-pressure-rectangle}
\end{align}
In the entrywise order from
Proposition~\ref{prop:qc-closed-rectangle-matrix}, these bounds have the
matrix form
\begin{equation}
 \mathbf C_{N,M}(T)
 \preceq_{\rm ent}
 \left(
 \begin{array}{ccc}
  T^{1/2}\mathsf Y_{n,M+1}^{\rm sm}
  &T^{1/2}\mathsf Y_{n+1,M-1}^{\rm sm}
  &TA_M\displaystyle\sum_{a'=0}^{n}\binom n{a'}
    \mathsf Y_{a',M+1}^{\rm sm}
    \mathsf Y_{n-a',M+1}^{\rm sm}
 \end{array}
 \right)_{n=0}^{N}.
 \label{eq:app-small-rectangle-matrix}
\end{equation}
Every right-hand side is an explicit finite function of \(\nu,T,N,M\) and
the datum norms
\begin{equation}
 \norm a_2,\qquad
 \norm{\Lambda^{1/2}a}_2,\qquad
 \left\{\norm{\Lambda^ra}_2:
  1\le r\le M+2N+1\right\}.
 \label{eq:app-small-data-derivative-ceiling}
\end{equation}
Indeed, induction in \(n\) shows that
\(\mathsf Y_{n,m}^{\rm sm}\) uses only the entries
\(U_j^{\rm sm}\) with \(-1\le j\le m+2n\); both velocity columns of
\eqref{eq:app-small-rectangle-matrix} consequently stop at
\(j=M+2N+1\), and its pressure column stops no later.
\end{proposition}

\begin{proof}
Equation \eqref{eq:app-small-data-Hm} proves the base
\eqref{eq:app-small-time-row} for \(m\ge1\); the kinetic-energy identity
proves it for \(m=0\), and \(H^0\hookrightarrow H^{-1}\) proves it for
\(m=-1\).  Apply \(\Pp\) to the temporal recurrence.
This is the \(\alpha=0\) temporal block-row of
\eqref{eq:formal-mixed-projected-matrix-row}; its factorial conjugate is
\eqref{eq:formal-factorial-velocity-convolution}.
Equation~\eqref{eq:qc-product-constant} at order \(m+1\) and induction give
\eqref{eq:app-small-time-row} through
\eqref{eq:app-small-Yrec}; for \(m=-1\), use the order-zero product
estimate and \(H^0\hookrightarrow H^{-1}\).  Integration in time gives
\eqref{eq:app-small-rectangle}.  The pressure recurrence and
\eqref{eq:qc-pressure-product-constant} at order \(M\), followed by
integration and summation in \(n\), give
\eqref{eq:app-small-pressure-rectangle}.  The same rowwise bounds give
\eqref{eq:app-small-rectangle-matrix}; taking the Frobenius norm of its first
two columns and summing its pressure column recovers
\eqref{eq:app-small-rectangle} and
\eqref{eq:app-small-pressure-rectangle}.
\end{proof}

\subsection{Low-rectangle and scale-local bounds}

The following results use distinct outputs of the compatible rectangles:
mixed space--time embeddings, pressure completion, and integrable Lipschitz
action.

\begin{proposition}[Three low-rectangle regularity classes]
\label{prop:app-low-rectangle-classes}
Let \(u\) be the global solution generated by
Theorem~\ref{thm:ier-global-smoothness}, and let \(0<T<\infty\).  Then
\begin{align}
 \norm u_{L^\infty(0,T;L^3)}
 &\le S_{1,0,3}\mathfrak Z_{0,1}[u_0,\nu;T],
 \label{eq:app-LinfL3}\\
 \norm u_{L^2(0,T;L^\infty)}
 &\le S_{2,0,\infty}\mathfrak X_{0,2}[u_0,\nu;T],
 \label{eq:app-L2Linfty}\\
 \norm{\nabla u}_{L^2(0,T;L^3)}
 &\le S_{2,1,3}\mathfrak X_{0,2}[u_0,\nu;T].
 \label{eq:app-grad-L2L3}
\end{align}
In particular, each of the three norms is finite on every finite horizon,
with the stated datum, viscosity, and horizon dependence.
\end{proposition}

\begin{proof}
Use \eqref{eq:qc-sharp-mixed-Lab} with
\((n,\alpha,a,b,q,k)=(0,0,\infty,3,1,0)\) to obtain
\eqref{eq:app-LinfL3}.  Use \eqref{eq:qc-sharp-mixed-L2b} with
\((n,\alpha,b,q,k)=(0,0,\infty,2,0)\) and
\((0,0,3,2,1)\) to obtain \eqref{eq:app-L2Linfty} and
\eqref{eq:app-grad-L2L3}.  The strict embedding conditions are,
respectively, \(1>1/2\), \(2>3/2\), and \(2>3/2\).
\end{proof}

\begin{corollary}[Uniform extinction of parabolic scale content]
\label{thm:app-scale-content-extinction}
Let \(\nu>0\), let \(u_0\in\HsSigma(\R^3)\), and let \((u,p)\) be
the pressure-normalized global pair supplied by
Theorem~\ref{thm:ier-global-smoothness}.  This is a post-regularity
consequence for the Caffarelli--Kohn--Nirenberg scale-normalized quantity
\cite{CKN1982}; it is not used as an \(\varepsilon\)-regularity input or
continuation criterion.
For \(0<T<\infty\), define
\begin{align}
 U_T&:=S_{2,0,\infty}\mathfrak Z_{0,2}[u_0,\nu;T],\notag\\
 P_T&:=S_{2,0,\infty}
 \mathfrak Q^{(\infty)}_{0,2}[u_0,\nu;T].
 \label{eq:app-scale-content-majorants}
\end{align}
Let \(z_0=(t_0,x_0)\), \(\rho>0\), and suppose that
\[
 Q_\rho(z_0):=(t_0-\rho^2,t_0)\times B_\rho(x_0)
 \subset(0,T)\times\R^3.
\]
Then the scale-normalized velocity--pressure content satisfies
\begin{equation}
 \rho^{-2}\iint_{Q_\rho(z_0)}
 \left(|u|^3+|p|^{3/2}\right)\,\dd x\dd t
 \le\frac{4\pi}{3}\rho^3
 \left(U_T^3+P_T^{3/2}\right).
 \label{eq:app-scale-content-bound}
\end{equation}
Consequently,
\begin{equation}
 \lim_{\rho\downarrow0}
 \sup_{\{z_0:\,Q_\rho(z_0)\subset(0,T)\times\R^3\}}
 \rho^{-2}\iint_{Q_\rho(z_0)}
 \left(|u|^3+|p|^{3/2}\right)\,\dd x\dd t=0.
 \label{eq:app-scale-content-uniform-limit}
\end{equation}
\end{corollary}

\begin{proof}
Equation~\eqref{eq:qc-sharp-mixed-sup}, followed by
\eqref{eq:qc-embedding-constant} with \((s,k,b)=(2,0,\infty)\), gives
\[
 \norm u_{L^\infty((0,T)\times\R^3)}\le U_T.
\]
Equations~\eqref{eq:qc-sharp-pressure-sup} and
\eqref{eq:qc-embedding-constant} give
\[
 \norm p_{L^\infty((0,T)\times\R^3)}\le P_T.
\]
Since \(|B_\rho|=\frac{4\pi}{3}\rho^3\) and the time length of
\(Q_\rho\) is \(\rho^2\),
\[
 \iint_{Q_\rho(z_0)}\left(|u|^3+|p|^{3/2}\right)\,\dd x\dd t
 \le\frac{4\pi}{3}\rho^5
 \left(U_T^3+P_T^{3/2}\right).
\]
Division by \(\rho^2\) proves \eqref{eq:app-scale-content-bound}; its
right-hand side is independent of the center and tends to zero with
\(\rho\), which proves \eqref{eq:app-scale-content-uniform-limit}.
\end{proof}

\subsection{Vorticity and material flow}

\begin{theorem}[Global vorticity amplification and dissipation]
\label{thm:app-vorticity-amplification}
Let \(\nu>0\), let \(u_0\in\HsSigma(\R^3)\), and let \((u,p)\) be the
pressure-normalized global pair supplied by
Theorem~\ref{thm:ier-global-smoothness}.  Put
\[
 \omega:=\nabla\times u,
 \qquad
 \omega_0:=\nabla\times u_0,
 \qquad
 L_T^\omega:=S_{3,1,\infty}T^{1/2}
 \mathfrak X_{0,3}[u_0,\nu;T].
\]
For every \(T<\infty\) and \(2\le r\le\infty\),
\begin{align}
 \sup_{0\le t\le T}\norm{\omega(t)}_{L^r}
 &\le e^{L_T^\omega}\norm{\omega_0}_{L^r},
 \label{eq:app-vorticity-Lr-amplification}\\
 \int_0^T\norm{(\omega\cdot\nabla)u(t)}_{L^r}\dd t
 &\le L_T^\omega e^{L_T^\omega}\norm{\omega_0}_{L^r}.
 \label{eq:app-vortex-stretching-action}
\end{align}
If \(2\le r<\infty\), then
\begin{equation}
 \int_0^T\norm{\nabla|\omega(t)|^{r/2}}_2^2\dd t
 \le \frac{r}{4\nu(r-1)}e^{rL_T^\omega}
       \norm{\omega_0}_{L^r}^r.
 \label{eq:app-vorticity-dissipation}
\end{equation}
\end{theorem}

\begin{proof}
Taking the curl of the momentum equation gives the three-dimensional
vorticity equation
\begin{equation}
 \partial_t\omega+(u\cdot\nabla)\omega
 = (\omega\cdot\nabla)u+\nu\Delta\omega.
 \label{eq:app-vorticity-equation}
\end{equation}
Fix \(2\le r<\infty\).  Pair
\eqref{eq:app-vorticity-equation} with
\(|\omega|^{r-2}\omega\).  To justify this power test, choose
\(\chi\in C_c^\infty(\R^3)\), with
\(0\le\chi\le1\), \(\chi=1\) on \(B_1\), and \(\chi=0\) outside
\(B_2\), put \(\chi_R(x)=\chi(x/R)\), and first test against
\[
 \chi_R\Phi_\varepsilon(\omega),
 \qquad
 \Phi_\varepsilon(z):=(|z|^2+\varepsilon)^{(r-2)/2}z.
\]
Define the primitive
\begin{equation}
 F_\varepsilon(z)
 :=\frac{(|z|^2+\varepsilon)^{r/2}-\varepsilon^{r/2}}{r}.
 \label{eq:app-vorticity-regularized-primitive}
\end{equation}
Then \(\nabla_zF_\varepsilon=\Phi_\varepsilon\), and incompressibility gives
\begin{equation}
 \int_{\R^3}\chi_R(u\cdot\nabla)\omega\cdot
 \Phi_\varepsilon(\omega)\dd x
 =-\int_{\R^3}F_\varepsilon(\omega)\,
 u\cdot\nabla\chi_R\dd x.
 \label{eq:app-vorticity-transport-cutoff}
\end{equation}
With \(A_R:=B_{2R}\setminus B_R\),
\begin{align}
 \left|\int F_\varepsilon(\omega)\,
 u\cdot\nabla\chi_R\dd x\right|
 &\le\frac{C_r}{R}\norm u_\infty
 \left(\norm\omega_{L^r(A_R)}^r
 +\varepsilon^{(r-2)/2}\norm\omega_{L^2(A_R)}^2\right),
 \label{eq:app-vorticity-transport-cutoff-bound}\\
 \left|\int \nabla\chi_R\cdot\nabla\omega\,
 \Phi_\varepsilon(\omega)\dd x\right|
 &\le\frac{C_r}{R}\norm{\nabla\omega}_{L^2(A_R)}
 \left(\norm\omega_{L^{2r-2}(A_R)}^{r-1}
 +\varepsilon^{(r-2)/2}\norm\omega_{L^2(A_R)}\right).
 \label{eq:app-vorticity-diffusion-cutoff-bound}
\end{align}
All norms on the right are finite because \(u,\omega\) belong to every
spatial Sobolev space.  Their annular tails vanish as \(R\to\infty\), so
\eqref{eq:app-vorticity-transport-cutoff} converges to the transport
cancellation and the diffusion cutoff term vanishes.

The remaining diffusion form is nonnegative:
\begin{align}
 \nabla\omega:D\Phi_\varepsilon(\omega)\nabla\omega
 ={}&(|\omega|^2+\varepsilon)^{(r-2)/2}|\nabla\omega|^2
 \notag\\
 &+(r-2)(|\omega|^2+\varepsilon)^{(r-4)/2}
 \sum_{j=1}^3(\omega\cdot\partial_j\omega)^2.
 \label{eq:app-vorticity-regularized-diffusion}
\end{align}
After the limit \(R\to\infty\), dominated convergence applies to
\(F_\varepsilon(\omega)\) and to the stretching term, while Fatou's lemma
applies to \eqref{eq:app-vorticity-regularized-diffusion} as
\(\varepsilon\downarrow0\).  This yields the unregularized chain rule and
diffusion inequality used below.  The stretching term
satisfies
\[
 \left|\int_{\R^3}(\omega\cdot\nabla)u\cdot
       |\omega|^{r-2}\omega\,\dd x\right|
 \le\norm{\nabla u}_\infty\norm\omega_r^r.
\]
For the diffusion term, spatial integration by parts gives
\begin{align*}
 -\int_{\R^3}\Delta\omega\cdot|\omega|^{r-2}\omega\,\dd x
 &\ge(r-1)\int_{\R^3}|\omega|^{r-2}|\nabla|\omega||^2\dd x\\
 &=\frac{4(r-1)}{r^2}
   \norm{\nabla|\omega|^{r/2}}_2^2.
\end{align*}
Hence
\begin{equation}
 \frac1r\frac{\dd}{\dd t}\norm\omega_r^r
 +\frac{4\nu(r-1)}{r^2}
   \norm{\nabla|\omega|^{r/2}}_2^2
 \le\norm{\nabla u}_\infty\norm\omega_r^r.
 \label{eq:app-vorticity-Lr-energy}
\end{equation}
Combining \eqref{eq:app-vorticity-Lr-energy} with the quantitative rectangle
estimate \eqref{eq:qc-sharp-lipschitz-action} gives
\[
 \int_0^T\norm{\nabla u(t)}_\infty\dd t\le L_T^\omega.
\]
Gronwall's inequality proves
\eqref{eq:app-vorticity-Lr-amplification} for finite \(r\).  Multiplying
\eqref{eq:app-vorticity-Lr-energy} by
\(
 \exp(-r\int_0^t\norm{\nabla u(s)}_\infty\dd s)
\)
and integrating proves \eqref{eq:app-vorticity-dissipation}.  Since
\(\omega_0,\omega(t)\in L^2\cap L^\infty\),
\(\norm f_r\to\norm f_\infty\) as \(r\to\infty\) for each such \(f\).
Consequently,
\[
 \sup_{0\le t\le T}\norm{\omega(t)}_\infty
 \le\liminf_{r\to\infty}
      \sup_{0\le t\le T}\norm{\omega(t)}_r
 \le e^{L_T^\omega}\lim_{r\to\infty}\norm{\omega_0}_r
 =e^{L_T^\omega}\norm{\omega_0}_\infty.
\]
This is \eqref{eq:app-vorticity-Lr-amplification} at the endpoint.  Finally,
\[
 \norm{(\omega\cdot\nabla)u(t)}_r
 \le\norm{\nabla u(t)}_\infty\norm{\omega(t)}_r;
\]
integration and \eqref{eq:app-vorticity-Lr-amplification} give
\eqref{eq:app-vortex-stretching-action}.
\end{proof}

\begin{theorem}[Quantitative volume-preserving material flow]
\label{thm:app-material-flow}
Let \(\nu>0\), let \(u_0\in\HsSigma(\R^3)\), and let \((u,p)\) be
the pressure-normalized global pair supplied by
Theorem~\ref{thm:ier-global-smoothness}.
For \(0<T<\infty\), put
\begin{align}
 D_T&:=S_{2,0,\infty}T^{1/2}
 \mathfrak X_{0,2}[u_0,\nu;T],\notag\\
 L_T&:=S_{3,1,\infty}T^{1/2}
 \mathfrak X_{0,3}[u_0,\nu;T].
 \label{eq:app-flow-actions}
\end{align}
For every \(a\in\R^3\), the problem
\begin{equation}
 \partial_tX(t,a)=u(t,X(t,a)),
 \qquad X(0,a)=a,
 \label{eq:app-flow-ode}
\end{equation}
has a unique solution on \([0,T]\).  For every \(t\in[0,T]\), the map
\(X(t,\cdot):\R^3\to\R^3\) is a smooth volume-preserving diffeomorphism
and
\begin{align}
 |X(t,a)-a|&\le D_T,
 \label{eq:app-flow-displacement}\\
 e^{-L_T}|a-b|
 &\le|X(t,a)-X(t,b)|
 \le e^{L_T}|a-b|,
 \label{eq:app-flow-bilipschitz}\\
 \det D_aX(t,a)&=1.
 \label{eq:app-flow-jacobian}
\end{align}
\end{theorem}

\begin{proof}
Equations~\eqref{eq:qc-sharp-L2Linfty} and
\eqref{eq:qc-sharp-lipschitz-action} give
\begin{align*}
 \int_0^T\norm{u(t)}_\infty\,\dd t&\le D_T,\\
 \int_0^T\norm{\nabla u(t)}_\infty\,\dd t&\le L_T.
\end{align*}
Picard iteration on successive compact time intervals gives local
existence and uniqueness for \eqref{eq:app-flow-ode}.  The first integral
bound prevents escape to spatial infinity and extends every trajectory to
\([0,T]\); it also proves \eqref{eq:app-flow-displacement}.  For two
trajectories, the second integral bound and Gronwall's inequality give the
upper bound in \eqref{eq:app-flow-bilipschitz}.  Solving the same equation
backward from time \(t\) produces the inverse map and gives the lower
bound.  Spatial differentiation of \eqref{eq:app-flow-ode} and induction
give smoothness of the map and its inverse.  Finally, Jacobi's formula and
\(\nabla\cdot u=0\) give
\[
 \partial_t\det D_aX
 =(\nabla\cdot u)(t,X)\det D_aX=0,
 \qquad \det D_aX(0,a)=1,
\]
which proves \eqref{eq:app-flow-jacobian}.
\end{proof}

\subsection{Relative energy and the Leray–Hopf class}

The weak class is stated in the whole-space form originating with Leray and
Hopf \cite{Leray1934,Hopf1951}.  The argument is the classical weak--strong
uniqueness mechanism associated with Serrin~\cite{Serrin1962}, written here
against the pressure-complete smooth history and derived directly from the
weak formulation and the energy inequality.

\begin{definition}[Three-dimensional Leray--Hopf solution]
\label{def:app-leray-hopf}
Let \(\nu>0\), \(T>0\), and \(v_0\in H^0_\sigma(\R^3)\).  A Leray--Hopf solution on
\([0,T]\) is a pair \((v,q)\) such that
\[
 v\in C_w([0,T];H^0_\sigma)
 \cap L^\infty(0,T;L^2_\sigma)
 \cap L^2(0,T;\dot H^1_\sigma),
 \qquad
 q\in\mathcal D'((0,T)\times\R^3),
\]
the equations
\[
 \partial_tv+(v\cdot\nabla)v+\nabla q=\nu\Delta v,
 \qquad
 \nabla\cdot v=0
\]
hold in distributions, \(v(t)\to v_0\) strongly in \(L^2\) as
\(t\downarrow0\), and
\begin{equation}
 \frac12\norm{v(t)}_2^2
 +\nu\int_0^t\norm{\nabla v(s)}_2^2\dd s
 \le\frac12\norm{v_0}_2^2
 \label{eq:app-leray-energy}
\end{equation}
for every \(t\in[0,T]\).  Every divergence-free test field
\(\phi\in C_c^\infty([0,T)\times\R^3;\R^3)\) satisfies
\begin{equation}
 \int_0^T\left[
 -\ip v{\partial_t\phi}
 -\int_{\R^3}(v\otimes v):\nabla\phi\dd x
 +\nu\ip{\nabla v}{\nabla\phi}
 \right]\dd t
 =\ip{v_0}{\phi(0)}.
 \label{eq:app-leray-weak-form}
\end{equation}
\end{definition}

A global Leray--Hopf solution at viscosity \(\nu\) is a pair on
\([0,\infty)\) whose restriction to every finite \([0,T]\) is a
Leray--Hopf solution in Definition~\ref{def:app-leray-hopf}, with the same
datum and viscosity.

\begin{lemma}[Solenoidal cutoff density and time commutation]
\label{lem:app-solenoidal-density}
Let \(\mathscr A\) be a finite subset of
\(\{(r,p):r\in\Nzero,\ 2\le p<\infty\}\).  There are linear operators
\(\mathcal T_k\), independent of time, such that, for every
\(f\in H^\infty_\sigma(\R^3)\),
\[
 \mathcal T_kf\in C_{c,\sigma}^\infty(\R^3)
\]
and, for every \((r,p)\in\mathscr A\),
\begin{align}
 \sup_k\norm{\mathcal T_kf}_{W^{r,p}}
 &\le C_{\mathscr A}\norm f_{W^{r,p}},
 \label{eq:app-solenoidal-density-uniform}\\
 \norm{\mathcal T_kf-f}_{W^{r,p}}&\longrightarrow0.
 \label{eq:app-solenoidal-density-strong}
\end{align}
The convergence is uniform for \(f\) in a compact subset of the retained
intersection of \(W^{r,p}\) spaces.  If
\(F\in C^\infty([0,T];H^\infty_\sigma)\), then
\begin{equation}
 \partial_t^a(\mathcal T_kF)=\mathcal T_k(\partial_t^aF)
 \qquad(a\in\Nzero),
 \label{eq:app-solenoidal-density-time}
\end{equation}
and the convergence in
\eqref{eq:app-solenoidal-density-strong} is uniform in \(t\) for every
fixed temporal derivative.
\end{lemma}

\begin{proof}
Choose \(\chi\in C_c^\infty(B_2)\) with
\(\chi=1\) on \(B_1\), and put
\(\chi_R(x)=\chi(x/R)\) and
\(A_R=B_{2R}\setminus\overline{B_R}\).  For a smooth solenoidal \(f\), set
\[
 g_R:=\nabla\chi_R\cdot f.
\]
It is supported in \(A_R\), and
\[
 \int_{A_R}g_R\,\dd x
 =\int_{\R^3}\nabla\cdot(\chi_Rf)\,\dd x=0.
\]
\begin{samepage}
We now construct the required compact anti-divergence explicitly.  Fix
\(\rho\in C_c^\infty((-3,3))\) with \(\int_{\R}\rho=1\).  If
\(h\in C_c^\infty(B_2)\) and \(\int h=0\), set
\[
 m_1(y_2,y_3):=\int_{\R} h(s,y_2,y_3)\,\dd s,
 \qquad
 m_2(y_3):=\int_{\R} m_1(s,y_3)\,\dd s
\]
and define \(\mathcal Gh=(G_1,G_2,G_3)\) by
\begin{align*}
 G_1(y)
 &:=\int_{-\infty}^{y_1}
 \bigl(h(s,y_2,y_3)-\rho(s)m_1(y_2,y_3)\bigr)\,\dd s,\\
 G_2(y)
 &:=\rho(y_1)\int_{-\infty}^{y_2}
 \bigl(m_1(s,y_3)-\rho(s)m_2(y_3)\bigr)\,\dd s,\\
 G_3(y)
 &:=\rho(y_1)\rho(y_2)
 \int_{-\infty}^{y_3}m_2(s)\,\dd s.
\end{align*}
\par
\end{samepage}
The three integrands have zero total integral in their integration
variables; for the last one this is \(\int h=0\).  Thus
\(\mathcal Gh\in C_c^\infty((-3,3)^3;\R^3)\), and direct differentiation
gives
\[
 \nabla\cdot\mathcal Gh
 =h-\rho(y_1)m_1
  +\rho(y_1)m_1-\rho(y_1)\rho(y_2)m_2
  +\rho(y_1)\rho(y_2)m_2
 =h.
\]
Integration over intervals of length at most six, H\"older's inequality,
Fubini's theorem, and differentiation under the integral give, for
\(a\in\Nzero\) and \(1<p<\infty\),
\begin{equation}
 \norm{\mathcal Gh}_{W^{a,p}((-3,3)^3)}
 \le C_{a,p,\rho}\norm h_{W^{a,p}(B_2)}.
 \label{eq:app-fixed-compact-antidivergence}
\end{equation}
For \(g\in C_c^\infty(B_{2R})\) of mean zero, put
\[
 (\mathcal G_Rg)(x)
 :=R\bigl[\mathcal G(g(R\,\cdot))\bigr](x/R).
\]
Then \(\nabla\cdot\mathcal G_Rg=g\), its support lies in
\((-3R,3R)^3\), and scaling \eqref{eq:app-fixed-compact-antidivergence}
gives
\begin{align}
 \norm{\mathcal G_Rg}_{L^p}
 &\le C_pR\norm g_{L^p},\notag\\
 \norm{\nabla^j\mathcal G_Rg}_{L^p}
 &\le C_{j,p}\sum_{b=0}^{j}
 R^{1-j+b}\norm{\nabla^bg}_{L^p},
 \qquad j\ge1.
 \label{eq:app-scaled-compact-antidivergence}
\end{align}
\begin{samepage}
Define
\[
 \mathcal S_Rf:=\chi_Rf-\mathcal G_Rg_R.
\]
Then \(\mathcal S_Rf\in C_{c,\sigma}^\infty(\R^3)\), with support in
\((-3R,3R)^3\).
\end{samepage}
Leibniz' rule gives, for \(0\le a\le r-1\),
\[
 \norm{\nabla^ag_R}_{L^p(A_R)}
 \le C_a\sum_{b=0}^{a}
 R^{-1-a+b}\norm{\nabla^bf}_{L^p(A_R)}.
\]
Combining this estimate with
\eqref{eq:app-scaled-compact-antidivergence} and the analogous cutoff
estimate yields
\begin{equation}
 \norm{\mathcal S_Rf-f}_{W^{r,p}}
 \le C_{r,p}\sum_{b=0}^{r}
 (1+R^{b-r})\norm{\nabla^bf}_{L^p(\R^3\setminus B_R)}.
 \label{eq:app-solenoidal-cutoff-tail}
\end{equation}
For \(r=0\), the same conclusion follows from
\eqref{eq:app-scaled-compact-antidivergence} and
\(\norm{g_R}_p\le C R^{-1}\norm f_{L^p(A_R)}\).
The right side tends to zero as \(R\to\infty\), and the same calculations
give operator norms bounded independently of \(R\).

Choose \(R_k\uparrow\infty\) and set
\(\mathcal T_k:=\mathcal S_{R_k}\).  Estimate
\eqref{eq:app-solenoidal-cutoff-tail} holds simultaneously for the finite
family \(\mathscr A\), while the preceding operator estimates give
\eqref{eq:app-solenoidal-density-uniform} simultaneously for that family.
Strong convergence plus the uniform operator bound makes convergence
uniform on compact subsets: cover a compact set by finitely many norm balls
and use the bound on each remainder.  Each \(\mathcal T_k\) is independent
of time, so differentiation commutes with it.  The ranges of
\(\partial_t^aF\) on \([0,T]\) are compact in every retained space, which
proves the final assertion.
\end{proof}

\begin{lemma}[Admissible strong-history test]
\label{lem:app-cross-energy}
Let \(u\) be the global solution from
Theorem~\ref{thm:ier-global-smoothness}, and let \(v\) be a Leray--Hopf
solution on \([0,T]\) at the same viscosity.  Put \(w:=v-u\).  Then, for every
\(t\in[0,T]\),
\begin{align}
 \ip{v(t)}{u(t)}
 &=
 \ip{v_0}{u_0}
 +\int_0^t\int_{\R^3}
       (w\cdot\nabla u)\cdot w\dd x\dd s
 \notag\\
 &\quad
 -2\nu\int_0^t\ip{\nabla v}{\nabla u}\dd s.
 \label{eq:app-cross-energy}
\end{align}
\end{lemma}

\begin{proof}
The energy class and spatial interpolation give
\[
 v\in L^4(0,T;L^3),
 \qquad
 \norm v_{L^4L^3}^4
 \le C\norm v_{L^\infty L^2}^2
       \norm{\nabla v}_{L^2L^2}^2,
\]
so \(v\otimes v\in L^2(0,T;L^{3/2})\).

Fix \(0<t<T\).  Solenoidal density on \(\R^3\), applied uniformly to the
compact sets \(u([0,T])\) and \(u_t([0,T])\) in their finite Sobolev spaces,
gives
\[
 u^{(k)}\in C^\infty([0,T];C_{c,\sigma}^\infty(\R^3))
\]
such that
\begin{equation}
 u^{(k)}\longrightarrow u
 \quad\text{in}\quad
 W^{1,1}(0,T;L^2)
 \cap L^2(0,T;H^1)
 \cap L^2(0,T;W^{1,3})
 \cap C([0,T];L^2).
 \label{eq:app-strong-test-density}
\end{equation}
Apply Lemma~\ref{lem:app-solenoidal-density} with
\[
 \mathscr A=\{(0,2),(1,2),(1,3)\}
\]
and set \(u^{(k)}=\mathcal T_ku\).  Its uniform compact-range conclusion
for \(u\) and \(u_t\), together with time commutation
\eqref{eq:app-solenoidal-density-time}, gives
\eqref{eq:app-strong-test-density}.

\begin{samepage}
For \(0<h<T-t\), choose \(\eta_h\in C_c^\infty([0,T))\) with
\(0\le\eta_h\le1\), \(\eta_h=1\) on \([0,t]\),
\(\eta_h=0\) on \([t+h,T]\), and \(\eta_h'\le0\).  The field
\(\eta_hu^{(k)}\) is an admissible test in
\eqref{eq:app-leray-weak-form}.  First let \(k\to\infty\).
The convergence in \eqref{eq:app-strong-test-density}, together with
\(v\otimes v\in L^2L^{3/2}\), permits passage in every term.
\par
\end{samepage}
\begin{samepage}
Choose the time cutoffs so that \(-\eta_h'\) is an approximate identity from the right
at \(t\), and let \(h\downarrow0\).  Weak continuity of \(v\), strong
continuity of \(u\), and the Lebesgue differentiation theorem give
\begin{align*}
 \ip{v(t)}{u(t)}
 &=
 \ip{v_0}{u_0}
 +\int_0^t\ip v{u_t}\dd s
 +\int_0^t\int_{\R^3}(v\otimes v):\nabla u\dd x\dd s\\
 &\quad-\nu\int_0^t\ip{\nabla v}{\nabla u}\dd s.
\end{align*}
\par
\end{samepage}
This identity first holds for almost every \(t\).  Its right-hand side is
continuous in \(t\), while \(t\mapsto\ip{v(t)}{u(t)}\) is continuous by
weak continuity of \(v\) and norm continuity of \(u\).  Hence it holds for
every \(0\le t\le T\), including \(t=T\) by passage from below.
\begin{samepage}
Using
\[
 u_t=-(u\cdot\nabla)u-\nabla p+\nu\Delta u
\]
is admissible because \(v\) is divergence free.  It gives
\[
 \int_0^t\ip v{u_t}\dd s
 =
 -\int_0^t\int_{\R^3}v\cdot(u\cdot\nabla)u\dd x\dd s
 -\nu\int_0^t\ip{\nabla v}{\nabla u}\dd s.
\]
Finally,
\[
 \int_{\R^3}
 \left[(v\otimes v):\nabla u-v\cdot(u\cdot\nabla)u\right]\dd x
 =
 \int_{\R^3}(w\cdot\nabla u)\cdot w\dd x,
\]
because
\[
 \int_{\R^3}(w\cdot\nabla u)\cdot u\dd x
 =\frac12\int_{\R^3}w\cdot\nabla|u|^2\dd x=0.
\]
This proves \eqref{eq:app-cross-energy}.
\end{samepage}
\end{proof}

\begin{theorem}[Quantitative Leray--Hopf stability]
\label{thm:app-relative-energy}
Let \(\nu>0\) and \(u_0\in\HsSigma(\R^3)\), and let \(u\) be the global strong history
from Theorem~\ref{thm:ier-global-smoothness}.  Fix \(T>0\), and let \(v\) be
any Leray--Hopf solution on \([0,T]\) at the same viscosity \(\nu\), with
datum \(v_0\in H^0_\sigma(\R^3)\).  Then, for \(0\le t\le T\),
\begin{align}
 &\frac12\norm{v(t)-u(t)}_2^2
 +\nu\int_0^t\norm{\nabla(v-u)(s)}_2^2\dd s
 \notag\\
 &\quad\le
 \frac12\norm{v_0-u_0}_2^2
 \exp\!\left(
 2\int_0^t\norm{\nabla u(s)}_\infty\dd s
 \right)
 \label{eq:app-relative-energy-action}\\
 &\quad\le
 \frac12\norm{v_0-u_0}_2^2
 \exp\!\left(
 2S_{3,1,\infty}T^{1/2}
 \mathfrak P_{0,2}[u_0,\nu;T]^{1/2}
 \right).
 \label{eq:app-relative-energy-producer}
\end{align}
The lower spatial rectangle also gives the viscosity-weighted estimate
\begin{align}
 &\frac12\norm{v(t)-u(t)}_2^2
 +\frac\nu2\int_0^t\norm{\nabla(v-u)(s)}_2^2\dd s
 \notag\\
 &\quad\le
 \frac12\norm{v_0-u_0}_2^2
 \exp\!\left(
 \frac{S_{2,0,\infty}^2}{\nu}
 \mathfrak X_{0,2}[u_0,\nu;T]^2
 \right).
 \label{eq:app-relative-energy-low-rectangle}
\end{align}
\end{theorem}

\begin{proof}
Combine \eqref{eq:app-leray-energy}, the strong energy equality
\eqref{eq:ier-global-energy}, and \eqref{eq:app-cross-energy}.  With
\(w=v-u\), this gives
\begin{equation}
 \frac12\norm{w(t)}_2^2
 +\nu\int_0^t\norm{\nabla w(s)}_2^2\dd s
 \le
 \frac12\norm{w(0)}_2^2
 -\int_0^t\int_{\R^3}(w\cdot\nabla u)\cdot w\dd x\dd s.
 \label{eq:app-relative-energy-raw}
\end{equation}
The integral term in \eqref{eq:app-relative-energy-raw} has absolute value at
most
\[
 \norm{\nabla u(s)}_\infty\norm{w(s)}_2^2.
\]
Gronwall's inequality bounds \(\norm{w(t)}_2^2\).  Substitution of that
bound back into \eqref{eq:app-relative-energy-raw} proves
\eqref{eq:app-relative-energy-action}.
\begin{samepage}
Proposition
\ref{prop:qc-continuation-norms} gives
\[
 \int_0^T\norm{\nabla u(t)}_\infty\dd t
 \le
 S_{3,1,\infty}T^{1/2}
 \mathfrak P_{0,2}[u_0,\nu;T]^{1/2},
\]
which proves \eqref{eq:app-relative-energy-producer}.
\end{samepage}

For the lower-rectangle estimate, integration by parts and
\(\nabla\cdot w=0\) give
\[
 \left|\int_{\R^3}(w\cdot\nabla u)\cdot w\,\dd x\right|
 =\left|\int_{\R^3}(w\cdot\nabla w)\cdot u\,\dd x\right|
 \le\norm u_\infty\norm w_2\norm{\nabla w}_2.
\]
Young's inequality in \eqref{eq:app-relative-energy-raw} gives
\begin{align*}
 \frac12\norm{w(t)}_2^2
 +\frac\nu2\int_0^t\norm{\nabla w(s)}_2^2\,\dd s
 \le\frac12\norm{w(0)}_2^2
 +\frac1{2\nu}\int_0^t\norm{u(s)}_\infty^2
 \norm{w(s)}_2^2\,\dd s.
\end{align*}
Set
\[
 A(t):=\frac1\nu\int_0^t\norm{u(s)}_\infty^2\,\dd s.
\]
Gronwall's inequality gives
\(\norm{w(t)}_2^2\le\norm{w(0)}_2^2e^{A(t)}\).  Substitution in the
preceding inequality gives
\[
 \frac12\norm{w(t)}_2^2
 +\frac\nu2\int_0^t\norm{\nabla w(s)}_2^2\,\dd s
 \le\frac12\norm{w(0)}_2^2e^{A(t)}.
\]
Finally, \eqref{eq:qc-sharp-L2Linfty} yields
\[
 A(t)\le\frac{S_{2,0,\infty}^2}{\nu}
 \mathfrak X_{0,2}[u_0,\nu;T]^2,
\]
which proves \eqref{eq:app-relative-energy-low-rectangle}.
\end{proof}

\begin{corollary}[Explicit weak-class stability for small critical data]
\label{cor:app-small-weak-stability}
Under the hypotheses of Theorem~\ref{thm:app-small-critical-global}, let
\(v\) be a global Leray--Hopf solution at the same viscosity \(\nu\), with
datum \(v_0\in H^0_\sigma(\R^3)\).
Then
\begin{align}
 &\frac12\norm{v(t)-u(t)}_2^2
 +\nu\int_0^t\norm{\nabla(v-u)(s)}_2^2\dd s
 \notag\\
 &\qquad\le
 \frac12\norm{v_0-a}_2^2
 \exp\!\left(2S_{3,1,\infty}tU_3^{\rm sm}\right)
 \qquad(t\ge0).
 \label{eq:app-small-weak-stability}
\end{align}
\end{corollary}

\begin{proof}
Equation~\eqref{eq:app-small-data-Hm} gives
\(\sup_{s\ge0}\norm{u(s)}_{H^3}\le U_3^{\rm sm}\).  Hence
\[
 \int_0^t\norm{\nabla u(s)}_\infty\dd s
 \le S_{3,1,\infty}tU_3^{\rm sm}.
\]
The solution constructed in Theorem~\ref{thm:app-small-critical-global} and
the global history from Theorem~\ref{thm:ier-global-smoothness}, with datum
\(a\), both restrict to members of
\(\mathscr M_\nu^3(0,T;a)\) on every finite interval.  The at-most-one
assertion in Theorem~\ref{thm:local-restart} identifies the two
pressure-normalized histories.
Apply \eqref{eq:app-relative-energy-action}.
\end{proof}

\begin{corollary}[Uniqueness and energy equality in the Leray--Hopf class]
\label{cor:app-leray-uniqueness}
Let \(\nu>0\) and \(u_0\in\HsSigma(\R^3)\).  If \((v,q)\) is any global
three-dimensional Leray--Hopf solution at viscosity \(\nu\), with datum
\(u_0\), then \(v\) agrees on every finite interval with the velocity of the
global solution from Theorem~\ref{thm:ier-global-smoothness}.  Thus \(v\) has
the smooth
representative
\[
 u\in C^\infty([0,\infty)\times\R^3;\R^3),
\]
and it satisfies
\[
 \frac12\norm{u(t)}_2^2
 +\nu\int_0^t\norm{\nabla u(s)}_2^2\dd s
 =\frac12\norm{u_0}_2^2
 \qquad(t\ge0).
\]
Its distributional pressure has the form
\begin{equation}
 q=p+c(t)
 \quad\text{in }\mathcal D'((0,\infty)\times\R^3),
 \label{eq:app-weak-pressure-gauge}
\end{equation}
where \(p\) is the pressure in \eqref{eq:ier-global-pressure} and \(c\) is a
scalar distribution in time.  The canonical Riesz representative
\(R_iR_j(u_i u_j)\) equals \(p\).  If \(q\) is represented by a locally
integrable function with \(q(t,\cdot)\in C_0(\R^3)\) for almost every \(t\),
then \(q=p\) almost everywhere.
\end{corollary}

\begin{proof}
Set \(v_0=u_0\) in Theorem~\ref{thm:app-relative-energy}.  The right-hand
side of \eqref{eq:app-relative-energy-producer} is zero, so
\(v=u\) in \(C_w([0,T];L^2)\cap L^2(0,T;\dot H^1)\).  Since \(T\) is
 arbitrary, the identity holds globally.  Smoothness and energy equality
 follow from \eqref{eq:ier-global-class} and
 \eqref{eq:ier-global-energy}.  If \(q\) is a distributional pressure paired
 with this velocity and \(p\) is the pressure in
 \eqref{eq:ier-global-pressure}, subtraction of the two momentum equations
 gives
 \[
  \nabla(q-p)=0
  \quad\text{in }\mathcal D'((0,\infty)\times\R^3).
 \]
 A distribution on a product domain whose spatial gradient vanishes is the
 pullback of a scalar distribution in time: testing first against spatial test
 functions of integral zero and then against one fixed test function of
 integral one gives \eqref{eq:app-weak-pressure-gauge}.  The Riesz formula for
 \(p\) proves the canonical-representative assertion.  In the function-valued
 case, \(q(t,x)-p(t,x)=c(t)\) for almost every \((t,x)\); both spatial slices
 tend to zero as \(\abs x\to\infty\), so \(c(t)=0\) for almost every \(t\).
\end{proof}
